\chapter{Trabajo futuro}

Este capítulo describe las líneas de trabajo futuro derivadas de las limitaciones metodológicas, clínicas y de ingeniería documentadas en la Sección~7.5. Cada propuesta se vincula con una restricción observada durante el desarrollo experimental o en la verificación del prototipo web (Capítulos~4 a~7).

Para superar la evaluación sobre una única partición fija 5/7--1/7--1/7 (Sección~4.3), se propone ejecutar validación cruzada de $k$ pliegues agrupados mediante \texttt{StratifiedGroupKFold}~\cite{pedregosa2011sklearn}. Este esquema cuantificará la varianza del Macro-F1 entre distintas particiones y evitará reutilizar un mismo conjunto de prueba en múltiples etapas experimentales~\cite{kapoor2023leakage}. Complementariamente, el uso de bancos con identificadores de paciente e inspección de duplicados eliminará el solapamiento residual del 0,14\,\% detectado en HAM10000~\cite{abhishek2025quality}.

Frente a la restricción de operar únicamente con píxeles dermatoscópicos (Sección~7.5), una segunda línea consiste en incorporar metadatos clínicos como edad, sexo y localización anatómica~\cite{tschandl2018ham10000}. La fusión multimodal de estas variables con los mapas de características de MobileNetV3Large~\cite{howard2019mobilenetv3} aportaría contexto clínico para desambiguar lesiones melanocíticas y queratósicas~\cite{amaral2025esmo}. En paralelo, un estudio de ablación por etapa medirá el aporte individual del filtro de mediana, la remoción de vello~\cite{wang2024hair} y el ajuste CLAHE~\cite{zuiderveld1994clahe}.

Para atender el reducido tamaño de prueba en \texttt{df} (17 imágenes), \texttt{vasc} (20) y \texttt{akiec} (46) (Tabla~\ref{tab:particion}), se plantea ampliar el entrenamiento con repositorios multicéntricos mayores~\cite{codella2019isic2018}. De igual forma, dado que HAM10000 proviene de poblaciones de Austria y Australia con predominio de pieles claras~\cite{tschandl2018ham10000}, se requiere incorporar imágenes de fototipos oscuros subrepresentados~\cite{daneshjou2022disparities}. Como demostraron Daneshjou et al.~\cite{daneshjou2022disparities}, el ajuste con cohortes diversas permite medir y reducir las brechas de desempeño entre subgrupos demográficos.

Con respecto a la brecha frente a la meta inicial de 0,80 de Macro-F1 (0,647 en la Tabla~\ref{tab:modos-decision}), se propone reemplazar la inicialización en ImageNet~\cite{deng2009imagenet}. El ajuste fino puede partir de modelos fundacionales específicos del dominio como PanDerm, preentrenado con más de dos millones de imágenes dermatológicas de once instituciones~\cite{yan2025panderm}. Esta transferencia desde representaciones dermatológicas previas debe evaluarse junto con técnicas de cuantización para conservar un tamaño compatible con la ejecución en navegador~\cite{ning2024onnxweb}.

En relación con la pérdida de especificidad observada al reducir el umbral escalar de melanoma a $t = 0{,}11$ (Tabla~\ref{tab:modos-decision}), se plantea explorar reglas de decisión sensibles al costo clínico~\cite{barata2023reinforcement}. A diferencia del corte sobre una sola probabilidad calibrada~\cite{guo2017calibration}, el enfoque de aprendizaje por refuerzo de Barata et al.~\cite{barata2023reinforcement} asigna recompensas asimétricas según la severidad de cada confusión diagnóstica. Incorporar una matriz de costos clínicos en las siete clases buscaría reducir los 50 melanomas confundidos con \texttt{nv} y los 21 confundidos con \texttt{bkl} (Sección~7.2).

Dado que el sistema se encuentra en nivel de madurez TRL~4 con pruebas retrospectivas (Sección~2.5), el paso necesario hacia un entorno asistencial es una evaluación clínica inicial supervisada~\cite{vasey2022decideai}. Siguiendo la guía DECIDE-AI~\cite{vasey2022decideai}, dicho estudio debe analizar la interacción de médicos y dermatólogos con los modos equilibrado y de alta sensibilidad en flujos reales de trabajo. Esta validación prospectiva requerirá cohortes de instituciones independientes de HAM10000 e ISIC 2018 y confirmación histopatológica de las lesiones sospechosas~\cite{amaral2025esmo,marchetti2023proveai,heinlein2024multicenter}.

Para formalizar la documentación de futuras versiones del clasificador y su validación en nuevos centros, se propone estructurar los informes técnicos bajo la guía TRIPOD+AI~\cite{collins2024tripodai}. La lista de verificación de TRIPOD+AI establece criterios estandarizados para reportar la procedencia de los datos, el tratamiento de valores faltantes, la calibración~\cite{sambyal2023calibration,penso2024calibration} y los protocolos de evaluación~\cite{collins2024tripodai}. Auditar cada ítem de dicha guía facilitará la reproducibilidad independiente de los experimentos y la comparación entre iteraciones del modelo~\cite{collins2024tripodai,maierhein2024metrics}.

Finalmente, para subsanar el alcance parcial de la verificación automática con \texttt{axe-core} (Sección~6.6), se plantea realizar una auditoría manual de accesibilidad sobre el aplicativo web~\cite{w3c2023wcag22}. Esta evaluación debe incluir pruebas con lectores de pantalla, recorrido por teclado sin ratón y validación con usuarios en distintos dispositivos móviles y de escritorio~\cite{w3c2023wcag22,wang2025browser}. Solo la combinación de pruebas automáticas y verificación humana permite evaluar la totalidad de los criterios de éxito de nivel A y AA de WCAG~2.2~\cite{w3c2023wcag22}.
